\section{Introduction}
\label{sec:intro}

High-dimension, low-sample-size (HDLSS) data situations occur in many areas of modern
science such as genetic microarrays, medical imaging, text recognition, finance and
chemometrics. In such data situations, the dimension $d$ is much larger than the sample
size $N$, so that the conventional multivariate procedures based on the large sample
asymptotic theory do not work.

Suppose we have independent and $d$-variate two populations, $\pi_i$, $i=1,2$, having
an unknown mean vector $\mu_i$ and unknown covariance matrix $\Sigma_i\succeq O$.
We have independent and identically distributed (i.i.d.) observations,
$x_{ij}$, $j=1,\dots,n_i$, from each $\pi_i$. We assume $n_i\ge 1$
($i=1,2$); the case $n_1=n_2=1$ is treated explicitly in Section~\ref{sec:example}.
Let $x_0$ be an observation vector of an individual belonging to one of the two
populations. We assume $x_0$ and $x_{ij}$'s are independent. Let
$N=n_1+n_2$ and $\Delta=\|\mu_1-\mu_2\|^2$, where $\|\cdot\|$ denotes the Euclidean
norm. Let $e(i)$ denote the error rate of misclassifying an individual from $\pi_i$
into the other class. We say that a classifier holds the consistency property if
\begin{equation}
e(1)+e(2)\longrightarrow 0
\quad\text{as }d\to\infty.
\label{eq:1.1}
\end{equation}
In the HDLSS context, \citet{hall2005} and \citet{marron2007} considered distance
weighted classifiers. \citet{chan2009} and \citet{aoshima2011} considered
distance-based classifiers. In particular, \citet{aoshima2011} gave the
misclassification rate adjusted classifier for multiclass, high-dimensional data in
which misclassification rates are no more than specified thresholds. On the other
hand, \citet{aoshima2014,aoshima2019} considered geometric classifiers based on a
geometric representation of HDLSS data, and discussed asymptotic properties and
optimality of the classifiers under high-dimension, non-sparse settings.

In the field of machine learning, a typical method of classification is the support
vector machine (SVM). Since HDLSS data are mostly separable by a hyperplane, one
usually considers the hard-margin SVM in the HDLSS context. For the linear SVM,
\citet{hall2005}, \citet{chan2009} and \citet{qiao2015} showed that the
misclassification rates tend to zero as $d\to\infty$ under certain severe conditions.
\citet{nakayama2017} investigated asymptotic properties of the linear SVM for HDLSS
data. They showed that the linear SVM is heavily biased when $n_1\neq n_2$, and that
the bias causes the strong inconsistency in which the misclassification rate tends to
one. In order to overcome such difficulties, they proposed a bias-corrected linear
SVM (BC-SVM) and showed that it holds the consistency property even for imbalanced
data. \citet{nakayama2020} investigated asymptotic properties of the SVM with the
Gaussian kernel and gave a choice of the scale parameter involved in the kernel.
\citet{nakayama2021,nakayama2022} considered the SVM for high-dimensional imbalanced
data in more general frameworks.

\subsection{The geometric representation and its limitation}

The above studies on the SVM rely on the geometric representation of HDLSS data.
Let $\theta_i=\tr(\Sigma_i)$ and $\theta=(\theta_1+\theta_2)/2$. Under the condition
\begin{equation}
\frac{\tr(\Sigma_i^2)}{\theta_i^2}\to 0
\quad\text{as }d\to\infty,
\label{eq:1.2}
\end{equation}
together with mild moment conditions, \citet{hall2005} and \citet{ahn2007} showed that
\[
\frac{\|x_{ij}-x_{ik}\|^2}{\theta_i}=2+o_P(1)\quad(j\neq k),
\qquad
\frac{\|x_{1j}-x_{2k}\|^2}{\theta}
=\frac{\theta_1+\theta_2}{\theta}+\frac{\Delta}{\theta}+o_P(1).
\]
That is, after rescaling by $\sqrt{\theta}$, all the observations are located at the
vertices of a regular simplex whose edge lengths are determined only by $\theta_i$
and $\Delta$. The randomness of the data vanishes in the limit. This deterministic
structure is the reason why the solution of the hard-margin SVM can be written in a
closed form in the HDLSS context, and why the discriminant function is described only
by $\Delta$ and
\begin{equation}
\kappa=\frac{\tr(\Sigma_1)}{n_1}-\frac{\tr(\Sigma_2)}{n_2}.
\label{eq:1.3}
\end{equation}
The quantity $\kappa$ is the bias term that the BC-SVM subtracts.

However, the condition~\eqref{eq:1.2} requires that the largest eigenvalue
$\lambda_{i1}$ of $\Sigma_i$ is negligible compared with $\theta_i$, since
$\lambda_{i1}^2\le\tr(\Sigma_i^2)$. In actual high-dimensional data such as gene
expression data, on the other hand, the eigenvalue structure is often spiked in the
sense that a few eigenvalues are much larger than the others and carry a substantial
proportion of the total variance. \citet{yata2012,yata2013} investigated HDLSS
asymptotic properties of the principal component analysis (PCA) under such spiked
models and proposed the noise-reduction (NR) methodology, which gives consistent
estimation of the spiked eigenvalues and eigenvectors when the sample size grows.
\citet{jung2009} showed that the sample principal component directions are
inconsistent when the spiked eigenvalue is of the same order as the trace and the
sample size is fixed.

Thus a natural question arises: what happens to the SVM in HDLSS settings when the
covariance matrices have dominant eigenvalues? To the best of our knowledge, this
question has not been studied. Since the geometric representation is the foundation
of the existing theory, we cannot expect that the known results carry over.

\subsection{Contributions of this paper}

In this paper, we consider asymptotic properties of the SVM in HDLSS settings under
a spiked model in which
\begin{equation}
\frac{\lambda_{is}}{\theta}\longrightarrow c_{is}\in(0,1),
\quad s=1,\dots,m_i,
\label{eq:1.4}
\end{equation}
where $m_i$ is a fixed integer. Our contributions are summarized as follows.

\begin{enumerate}[label=(\roman*)]
\item We show that the geometric representation does not hold under~\eqref{eq:1.4}.
We show that the Gram matrix of the rescaled data converges in distribution to a
random matrix, and we give its explicit form (Theorem~\ref{thm:1}). Moreover, we
show that the limiting Gram matrix is realized as the Gram matrix of an explicit
family of vectors in $\R^{1+m_1+m_2}$ (Lemma~\ref{lem:6}), namely, the HDLSS data
converge to a random configuration in a finite-dimensional space together with $N$
mutually orthogonal noise axes. We call this result the limiting configuration
theorem.

\item We show that the solution of the hard-margin SVM converges in distribution to
the solution of the limiting problem (Theorem~\ref{thm:2}), and that the discriminant
function converges in distribution to a non-degenerate random variable
(Theorem~\ref{thm:3}). Consequently, we show that the misclassification rates do not
tend to zero, that is, the SVM does not hold the consistency property
(Corollary~\ref{cor:1}). We emphasize that this inconsistency is different from the
strong inconsistency of \citet{nakayama2017}, in which the misclassification rate
tends to one. Here the misclassification rate tends neither to zero nor to one.

\item We give an example in which the limiting misclassification rate is obtained in
a closed form (Theorem~\ref{thm:4}). We show that a spike degrades the performance of
the SVM even when it is orthogonal to the mean difference. This phenomenon cannot
be observed in the existing theory.

\item We show that the bias term should be modified from~\eqref{eq:1.3} to
$\kappa_*=\tr(\Sigma_{1*})/n_1-\tr(\Sigma_{2*})/n_2$, where $\Sigma_{i*}$ denotes
the non-spiked part of $\Sigma_i$ (Proposition~\ref{prop:1}). We show that the
BC-SVM overcorrects the bias under the spiked model.

\item In order to overcome such difficulties, we propose a spike-corrected SVM
(SC-SVM), which removes the uninformative spiked directions by the NR methodology
and applies the BC-SVM with the modified bias term. We show that the SC-SVM holds
the consistency property when $n_i\to\infty$ (Theorem~\ref{thm:5}). On the other
hand, we show that any projection-based procedure does not hold the consistency
property when $N$ is fixed (Proposition~\ref{prop:3}), so that the growth of the
sample size is essential.

\item Finally, we check the performance of the classifiers by numerical simulations.
\end{enumerate}

The rest of the paper is organized as follows. In Section~\ref{sec:notation}, we
give the notation and the assumptions. In Section~\ref{sec:inner}, we investigate
asymptotic behaviour of the inner products of HDLSS data under the spiked model.
In Section~\ref{sec:config}, we give the limiting configuration theorem. In
Section~\ref{sec:svm}, we show the convergence of the SVM. In
Section~\ref{sec:disc}, we give the limit of the discriminant function and the
inconsistency. In Section~\ref{sec:example}, we give the explicit example. In
Section~\ref{sec:bias}, we give the modified bias term. In
Section~\ref{sec:scsvm}, we propose the SC-SVM. In Section~\ref{sec:fixed}, we
discuss the case that $N$ is fixed. In Section~\ref{sec:sim}, we give numerical
simulations. In Section~\ref{sec:conclude}, we give concluding remarks. The
mathematical tools used in Sections~\ref{sec:svm} and~\ref{sec:disc} are summarized
in Appendix~\ref{app:A}.
